\bookchapter{Глава 29: Цифровое здоровье — как выстроить здоровые отношения со смартфоном}

\chapterintro{Знакомо ли вам: вы хотели посмотреть короткие видео всего 5 минут, а подняли голову — прошло уже два часа? Или лежите ночью в кровати со смартфоном, и чем дольше листаете, тем бодрее себя чувствуете, а в итоге засыпаете только в два часа ночи? Смартфон изменил нашу жизнь, но при этом тихо пожирает наше время и здоровье. Эта глава не для того, чтобы «бросить смартфон», а чтобы помочь вам вернуть контроль над ним.}

\booksection{Сначала посмотрите на своё «экранное время» — цифры отрезвляют}

Прежде чем что-то менять, разберитесь в текущей ситуации. В смартфоне есть функция, которую многие никогда не открывали, — «Экранное время» (iPhone) или «Цифровое благополучие» (Android). Она показывает, сколько времени вы проводите в телефоне каждый день и на какие приложения оно уходит.

\begin{itemize}
  \item \textbf{iPhone}: «Настройки» > «Экранное время» > «Посмотреть все действия». Вы увидите ежедневное число разблокировок, количество раз, когда вы взяли телефон в руки, и время использования каждого приложения.
  \item \textbf{Android}: «Настройки» > «Цифровое благополучие и родительский контроль». У Huawei это называется «Здоровое использование телефона», у Xiaomi — «Управление экранным временем». Функции аналогичны, можно увидеть ежедневную статистику использования.
\end{itemize}

\textbf{80\% людей, впервые увидевших эти данные, пугаются}. Мой друг после просмотра обнаружил, что проводит в TikTok более 4 часов в день — это 28 часов в неделю, что эквивалентно подработке на полставки. Этот «цифровой шок» — первый шаг к переменам.

\booksection{Установите «лимит использования» для приложений — поставьте себе ограничитель скорости}

Когда вы знаете, в чём проблема, можно точечно настроить ограничения. Не нужно сразу удалять приложения — достаточно установить для них «дневной лимит»:

\begin{enumerate}
  \item \textbf{iPhone}: «Настройки» > «Экранное время» > «Лимиты приложений» > «Добавить лимит». Выберите приложение для ограничения (например, TikTok) и установите 30 минут в день. По истечении времени иконка приложения станет серой, а для его открытия потребуется ввести пароль.
  \item \textbf{Android}: «Настройки» > «Цифровое благополучие» > «Таймеры приложений». Выберите приложение и установите дневной лимит — по его истечении иконка приложения станет серой.
  \item \textbf{Маленький совет}: установите в качестве пароля лимита комбинацию, которую вы редко используете и с трудом вспоминаете. Если пароль слишком лёгкий, вы не сможете удержаться от «продления» каждый раз, и лимит станет бесполезным.
\end{enumerate}

\booksection{Правильное использование режима «Не беспокоить» — не отключение уведомлений, а избавление от тревожности}

У многих смартфон выдаёт десятки уведомлений в день, и каждое из них обрывает ваше внимание. Вам кажется, что вы «просто взглянули на уведомление», но на самом деле для повторной концентрации после каждого перерыва требуется 15–20 минут. Режим «Не беспокоить» — это не способ «отключиться от мира», а способ вернуть контроль над своим вниманием.

\begin{itemize}
  \item \textbf{Запуск по расписанию}: настройте режим «Не беспокоить» на автоматическое включение с 22:00 до 7:00. Тогда во сне телефон не разбудит вас из-за спам-СМС.
  \item \textbf{Разрешите проход важным контактам}: в настройках «Не беспокоить» установите «Разрешить вызовы от» на «Избранные контакты» или «Повторные вызовы» (второй звонок в течение 3 минут зазвонит). Так вы не пропустите экстренный вызов от близких, но и не будете беспокоены рекламными звонками.
  \item \textbf{Режим фокусировки на работе}: «Режим фокусирования» в iPhone позволяет создавать пользовательские сценарии — например, в режиме «Работа» разрешены только уведомления рабочих приложений, а в режиме «Чтение» все уведомления беззвучны.
\end{itemize}

\booksection{Защита зрения — не приложение, а привычка: правило 20-20-20}

При длительном использовании смартфона сухость и напряжение глаз неизбежны. Режим защиты зрения (фильтр синего света) немного помогает, но по-настоящему эффективно изменить привычки:

\begin{tcbinfo}{Правило 20-20-20}
Каждые\textbf{20 минут}работы за экраном переводите взгляд на объект в\textbf{20 футах (около 6 метров)}от вас и держите его\textbf{20 секунд}. Это простое действие эффективно снимает усталость ресничной мышцы.

\end{tcbinfo}

\begin{itemize}
  \item \textbf{Включите режим защиты зрения}: «Настройки» > «Экран» > «Защита зрения» (или «Ночной режим»). После включения цветовая температура экрана становится теплее, уменьшается синий свет. Рекомендуется включать автоматически с заката до рассвета.
  \item \textbf{Настройте яркость экрана}: в тёмном помещении экран тоже должен быть тусклым. Многие в темноте выкручивают яркость на максимум — это одна из самых вредных привычек. Включите «Автояркость» и убедитесь, что она работает.
  \item \textbf{Держите дистанцию}: расстояние от телефона до глаз должно быть не менее 30 см. Многие любят подносить смартфон слишком близко — ресничная мышца постоянно напряжена, и близорукость быстро прогрессирует.
\end{itemize}

\booksection{Самопомощь для шейного отдела «поколения с опущенной головой»}

Когда вы смотрите в смартфон с опущенной головой, шея испытывает нагрузку в 5–6 раз больше, чем в нормальном вертикальном положении. Длительное наклонение головы приводит к выпрямлению шейного лордоза, остеофитам, головокружению и онемению рук. Вот несколько простых способов улучшить ситуацию:

\begin{itemize}
  \item \textbf{Поднимайте телефон выше}: старайтесь держать смартфон на уровне глаз, а не наклоняться к нему, лежащему на коленях. Руки могут устать, но ваша шея скажет «спасибо».
  \item \textbf{Каждые 30 минут разминайте шею}: медленно поверните голову влево до упора, задержите на 3 секунды, вернитесь в исходное положение; затем вправо, задержите на 3 секунды. Повторите 5 раз. Затем выполните упражнение «рисование знака» — подбородком «напишите» в воздухе знак «звезда» (форма китайского иероглифа «ми»).
  \item \textbf{Стойте у стены 3 минуты в день}: затылок, лопатки, ягодицы и пятки должны касаться стены. Ежедневные 3 минуты такого стояния эффективно исправляют сутулость и привычку опускать голову.
\end{itemize}

\booksection{Сон и смартфон — привычка, которую нужно изменить в первую очередь}

Синий свет экрана смартфона подавляет выработку мелатонина, из-за чего вы «чувствуете себя всё бодрее». Многие думают, что «вот-вот уснут», но на самом деле синий свет обманывает мозг, и после того, как вы откладываете телефон, уснуть становится сложнее. Вот рекомендации по улучшению сна:

\begin{itemize}
  \item \textbf{Не пользоваться смартфоном за 1 час до сна}: это самое эффективное и самое сложное правило. Если не получается, хотя бы включите «Ночной режим / Защиту зрения», чтобы снизить уровень синего света.
  \item \textbf{Уберите телефон подальше}: поставьте зарядку в другом конце спальни, а не на прикроватной тумбочке. Так вы не сможете удержаться от очередной проверки уведомлений, а когда зазвонит будильник, вам придётся встать, чтобы выключить его — двойная польза.
  \item \textbf{Включите «Режим сна»}: режим фокусирования «Сон» в iPhone автоматически блокирует экран и упрощает экран блокировки в часы сна; на Android у разных брендов есть аналогичные функции.
  \item \textbf{Используйте обычный будильник}: купите обычный будильник за несколько десятков рублей вместо использования телефона. Вы удивитесь, как исчезнет соблазн «глянуть на время и заодно полистать ленту».
\end{itemize}

\booksection{Информационная тревожность — учитесь «выборочной слепоте»}

Объём информации, которую мы получаем ежедневно, равен тому, что древний человек получал за всю жизнь. Но человеческий мозг не приспособлен к обработке такого потока — перегрузка ведёт к тревожности, рассеянности и усталости от принятия решений. Решение — не «получать больше информации», а «активно её отсеивать»:

\begin{itemize}
  \item \textbf{Отписывайтесь}: откройте список всех публичных аккаунтов, блогеров и авторов, на которых подписаны, и задайте себе вопрос: «Получил ли я от них реальную пользу за последний месяц?» Если нет — отписывайтесь. Лучше меньше каналов, но качественных.
  \item \textbf{Отключите ненужные уведомления}: кроме WeChat, звонков, SMS и календаря, отключите уведомления у всех остальных приложений. Новости не нуждаются в push-уведомлениях — когда вы сами захотите их прочитать, вы их увидите.
  \item \textbf{Введите «информационный пост»}: выберите один день в неделю (например, воскресенье) и полностью откажитесь от новостей и коротких видео. Вы удивитесь, что мир не рухнет, а в понедельник будете чувствовать себя заметно бодрее.
\end{itemize}

\begin{tcbpractice}{Практика: недельный челлендж цифрового здоровья}
\begin{enumerate}
  \item Сегодня: откройте «Экранное время» или «Цифровое благополучие» и посмотрите статистику использования за прошлую неделю.
  \item Завтра: установите дневной лимит для самого «залипательного» развлекательного приложения (рекомендуется 30–60 минут).
  \item Послезавтра: настройте автоматическое включение режима «Не беспокоить» (22:00–07:00), разрешив звонки от избранных контактов.
  \item Четвёртый день: не пользоваться телефоном за 1 час до сна, поставить зарядку в другой конец спальни.
  \item Пятый день: отписаться от 5 публичных аккаунтов или блогеров, которые вы больше не читаете.
  \item Шестой день: попробуйте провести целый день в «информационном посте» — без новостей и коротких видео.
  \item Седьмой день: посмотрите отчёт об экранном времени за эту неделю и сравните с прошлой.
\end{enumerate}

\end{tcbpractice}

\begin{tcbsummary}{Итоги главы}
\begin{itemize}
  \item Сначала посмотрите данные об экранном времени — осознание текущего состояния это первый шаг к изменениям
  \item Установите лимиты для развлекательных приложений, настройте расписание «Не беспокоить» и верните контроль над вниманием
  \item Правило 20-20-20 — самая простая и эффективная привычка для защиты зрения
  \item Когда вы смотрите в телефон с опущенной головой, шея испытывает нагрузку в 5–6 раз больше — поднимайте телефон и регулярно разминайте шею
  \item Не смотрите в телефон за час до сна, ставьте зарядку в место, куда нужно тянуться
  \item Лекарство от информационной тревожности — «активное экранирование», а не «получать больше»
\end{itemize}

\end{tcbsummary}
